\section*{अध्याय \ref{ch:b3:solid-state} --- ठोस-अवस्था रसायन: बैंड, दोष और अर्धचालक}

\begin{solution}{exo:b3:solid-state:1}
$(E - \alpha)/|\beta| = -2\cos(k\pi/7)$: $-1.802$, $-1.247$, $-0.445$, $+0.445$, $+1.247$, $+1.802$। वे
$3.604|\beta|$ घेरते हैं, जो अनंत शृंखला की बैंड चौड़ाई $4|\beta|$ का पहले ही 90\,\% है।
\end{solution}

\begin{solution}{exo:b3:solid-state:2}
$2N$ ($N$ कक्षक, प्रत्येक में दो इलेक्ट्रॉन)। सोडियम, प्रति परमाणु एक 3s इलेक्ट्रॉन के साथ, अपना s बैंड
आधा भरता है: एक धातु। मैग्नीशियम में दो हैं और वह s बैंड को ठीक पूरा भर देता,
परंतु 3s और 3p बैंड अतिव्यापित होते हैं, अतः सबसे ऊँचे अधिकृत स्तर एक
आंशिक रूप से भरे संयुक्त बैंड के हैं: वह भी एक धातु।
\end{solution}

\begin{solution}{exo:b3:solid-state:3}
$\lambda \approx \qty{1239.8}{eV.nm}/E_g$: गैलियम फ़ॉस्फ़ाइड \qty{549}{nm} (हरा); गैलियम नाइट्राइड \qty{365}{nm}
(निकट पराबैंगनी)। नीले डायोड इंडियम के साथ मिश्रधातुकृत गैलियम नाइट्राइड का उपयोग करते हैं, जिसका
अंतराल छोटा है।
\end{solution}

\begin{solution}{exo:b3:solid-state:4}
$\ce{CaCl2} \xrightarrow{\ce{NaCl}} \mathrm{Ca_{Na}^{\bullet} + V_{Na}' + 2\,Cl_{Cl}^{\times}}$: दो धनायन और दो ऋणायन स्थल, जैसे NaCl में; एक Ca और
दो Cl; आवेश $+1 - 1 = 0$। $\ce{Y2O3} \xrightarrow{\ce{ZrO2}} \mathrm{2\,Y_{Zr}' + 3\,O_O^{\times} + V_O^{\bullet\bullet}}$: दो धनायन और चार ऑक्सीजन स्थल; दो Y
और तीन O; आवेश $-2 + 2 = 0$।
\end{solution}

\begin{solution}{exo:b3:solid-state:5}
\qty{300}{K} पर $T^{3/2} = 5196$, $kT = \qty{0.02585}{eV}$। सिलिकॉन: $\sqrt{N_cN_v} = \qty{2.42e19}{cm^{-3}}$, $\eu^{-1.125/0.05170} = \num{3.55e-10}$,
$n_i = \qty{8.6e9}{cm^{-3}}$। जर्मेनियम: $\sqrt{N_cN_v} = \qty{7.16e18}{cm^{-3}}$, $\eu^{-0.661/0.05170} = \num{2.80e-6}$, $n_i = \qty{2.0e13}{cm^{-3}}$। अनुपात
$\num{4.3e-4}$: जर्मेनियम का छोटा अंतराल उसे दो हज़ार गुना से अधिक
वाहक देता है।
\end{solution}

\begin{solution}{exo:b3:solid-state:6}
$n = N_D = \qty{1.0e16}{cm^{-3}}$; $p = n_i^2/n = \num{1.0e20}/\num{1.0e16} = \qty{1.0e4}{cm^{-3}}$।
\end{solution}

\begin{solution}{exo:b3:solid-state:7}
$E_F - E_i = kT\ln(n/n_i) = \qty{0.02585}{eV} \times \ln(10^6) = \qty{0.36}{eV}$।
\end{solution}

\begin{solution}{exo:b3:solid-state:8}
\qty{800}{K} पर $kT = \qty{0.06894}{eV}$; $n/N = \eu^{-2.3/0.1379} = \num{5.7e-8}$; एक मोल में $\num{6.02e23} \times \num{5.7e-8} = \num{3.4e16}$
युग्म।
\end{solution}

\begin{solution}{exo:b3:solid-state:9}
एक कोष्ठिका का द्रव्यमान $\rho a^3 = \qty{5.70}{g.cm^{-3}} \times (\qty{4.300e-8}{cm})^3 = \qty{4.532e-22}{g}$ है, अर्थात्
$\qty{4.532e-22}{g} \times N_A = \qty{272.9}{g}$ प्रति मोल कोष्ठिकाएँ, या प्रति ऑक्सीजन \qty{68.23}{g}। तब
$(1 - x) \times 55.8 + 16.0 = 68.23$ से $1 - x = 0.936$ मिलता है: $x = 0.064$, $\mathrm{Fe_{0.936}O}$।
\end{solution}

\begin{solution}{exo:b3:solid-state:10}
$\ln(\sigma T)$: $-3.51$, $-1.30$, $0.36$, $1/T = 3.333$, $2.857$, $\qty{2.500e-3}{K^{-1}}$ पर; छोर-बिंदुओं के बीच ढलान
$(0.36 + 3.51)/(-\qty{8.33e-4}{K^{-1}}) = -\qty{4.65e3}{K}$ (मध्य बिंदु रेखा पर है), अतः
$E_a = \qty{4.65e3}{K} \times \qty{8.617e-5}{eV.K^{-1}} = \qty{0.40}{eV}$।
\end{solution}

\begin{solution}{exo:b3:solid-state:11}
$\tfrac12\ce{O2} \rightleftharpoons \mathrm{O_O^{\times} + V_M' + h^{\bullet}}$, $K = [\mathrm{V_M'}][\mathrm h^\bullet]/p(\ce{O2})^{1/2}$। आवेश संतुलन $[\mathrm h^\bullet] = [\mathrm{V_M'}]$, अतः
$[\mathrm h^\bullet]^2 = K\,p(\ce{O2})^{1/2}$ और $[\mathrm h^\bullet] \propto p(\ce{O2})^{1/4}$; चालकता \omterm{def:b3:solid-state:carriers}{कोटरों} का अनुसरण करती है।
\end{solution}

\begin{solution}{exo:b3:solid-state:12}
$n/N = \sqrt{N_i/N}\,\eu^{-\Delta H_F/2kT} = \sqrt 2\,\eu^{-1.4/(2 \times 0.05170)} = 1.414 \times \num{1.32e-6} = \num{1.9e-6}$।
\end{solution}

\begin{solution}{pb:b3:solid-state:1}
\textbf{1.} $\ce{Y2O3} \xrightarrow{\ce{ZrO2}} \mathrm{2\,Y_{Zr}' + 3\,O_O^{\times} + V_O^{\bullet\bullet}}$।

\textbf{2.} स्थल: दो धनायन स्थल और चार ऑक्सीजन स्थल (तीन भरे, एक रिक्त),
$\ce{ZrO2}$ का 1\,:\,2 अनुपात। द्रव्यमान: प्रत्येक ओर 2 Y और 3 O। आवेश: $2 \times (-1) + 2 = 0$।

\textbf{3.} दो यिट्रियम आयनों पर एक रिक्ति: प्रति यिट्रियम आधी।

\emergencystretch=3em \textbf{4.} प्रति $(\ce{ZrO2})_{0.92}(\ce{Y2O3})_{0.08}$ में 0.92 Zr और 0.16 Y हैं, 1.08 धनायन: $x = 0.16/1.08 = 0.148$,
$\mathrm{Zr_{0.852}Y_{0.148}O_{1.926}}$।

\textbf{5.} $(x/2)/2 = 0.0370$: 3.7\,\% ऑक्सीजन स्थल रिक्त हैं।

\textbf{6.} प्रति कोष्ठिका $8 \times 0.0370 = 0.296$ रिक्तियाँ।

\textbf{7.} $a^3 = (\qty{5.14e-8}{cm})^3 = \qty{1.358e-22}{cm^3}$: $0.296/\qty{1.358e-22}{cm^3} = \qty{2.2e21}{cm^{-3}}$।

\textbf{8.} प्रत्येक ऑक्साइड आयन को किसी पड़ोसी रिक्ति में \qty{1.0}{eV} के अवरोध के ऊपर से
कूदना होता है; सफल प्रयासों का अंश, $\eu^{-E_a/kT}$, $T$ के साथ बहुत तेज़ी से बढ़ता है।

\textbf{9.} \qty{1073.15}{K} पर $kT = \qty{0.09248}{eV}$: $A = \sigma T\eu^{E_a/kT} = 0.030 \times 1073.15 \times \eu^{10.81} = \qty{1.6e6}{S.K.cm^{-1}}$।

\textbf{10.} \qty{973.15}{K} पर: $\sigma = (A/T)\eu^{-11.92} = \qty{0.011}{S.cm^{-1}}$।

\textbf{11.} \qty{573.15}{K} पर: $\sigma = (A/T)\eu^{-20.25} = \qty{4.5e-6}{S.cm^{-1}}$।

\textbf{12.} $R = L/(\sigma S) = \qty{0.10}{cm}/(\sigma \times \qty{0.20}{cm^2})$: \qty{700}{\celsius} पर \qty{46}{\ohm}, \qty{300}{\celsius} पर $\qty{1.1e5}{\ohm}$।

\textbf{13.} $R < \qty{1}{k\ohm}$ के लिए $\sigma > \qty{5.0e-4}{S.cm^{-1}}$ चाहिए; $(A/T)\eu^{-E_a/kT} = \num{5.0e-4}$ को (प्रयत्न द्वारा) हल करने पर
$T \approx \qty{760}{K}$ मिलता है, लगभग \qty{490}{\celsius}।

\textbf{14.} स्टार्ट के बाद और कम भार पर निकास ठंडा होता है; तापक
विद्युत-अपघट्य को सेकंडों में उसके कार्य ताप तक ले आता है, ताकि नियंत्रण
शुरुआत से ही कार्य करे, जब अधिकांश प्रदूषक उत्सर्जित होते हैं।

\textbf{15.} $\ce{O2 + 4e- -> 2O^2-}$; \omterm{def:b3:solid-state:kroger-vink}{क्रोगर--विंक संकेतन} में $\ce{O2} + \mathrm{2\,V_O^{\bullet\bullet} + 4\,e'} \longrightarrow \mathrm{2\,O_O^{\times}}$।

\textbf{16.} वायु वाले पक्ष पर ऑक्सीजन अपचयित होती है और निकास वाले पक्ष पर ऑक्साइड आयन फिर से
ऑक्सीजन में ऑक्सीकृत होते हैं: सेल $\ce{O2}$ को वायु से निकास की ओर स्थानांतरित करता है,
$\Delta_rG = RT\ln(p_{\text{निकास}}/p_{\text{वायु}})$ के साथ, प्रति मोल $\ce{O2}$ और चार इलेक्ट्रॉनों पर: $E = -\Delta_rG/4F = (RT/4F)\ln(p_{\text{वायु}}/p_{\text{निकास}})$।

\textbf{17.} $RT/4F = 8.3145 \times 973.15/(4 \times 96485) = \qty{0.02096}{V}$।

\textbf{18.} शून्य: दोनों दाब बराबर हैं।

\textbf{19.} प्रति दशक $(RT/4F)\ln 10 = \qty{48.3}{mV}$।

\textbf{20.} $\qty{0.02096}{V} \times \ln(0.2095/0.010) = \qty{0.02096}{V} \times 3.04 = \qty{0.064}{V}$।

\textbf{21.} $\qty{0.02096}{V} \times \ln(0.2095/\num{1.0e-20}) = \qty{0.02096}{V} \times 44.49 = \qty{0.93}{V}$।

\textbf{22.} समृद्ध निकास में कार्बन मोनोऑक्साइड, हाइड्रोजन और अदग्ध
हाइड्रोकार्बन होते हैं; प्लैटिनम इलेक्ट्रोड पर बची थोड़ी ऑक्सीजन उनसे अभिक्रिया करती है,
और $p(\ce{O2})$ साम्यों $\ce{2CO + O2 <=> 2CO2}$ और $\ce{2H2 + O2 <=> 2H2O}$ से तय होता है, जो \qty{700}{\celsius} पर एक अत्यंत छोटा
मान छोड़ते हैं।

\textbf{23.} $p = 0.2095\,\eu^{-0.45/0.02096} = \qty{1.0e-10}{bar}$।

\textbf{24.} रससमीकरणमितीय अनुपात के आसपास निकास ऑक्सीजन के थोड़े आधिक्य से
CO और हाइड्रोजन के थोड़े आधिक्य में बदल जाता है: ईंधन में छोटे परिवर्तन के लिए $p(\ce{O2})$ दस की लगभग अठारह
घातों से गिरता है, और वोल्टता
\qty{0.1}{V} से नीचे से उछलकर लगभग \qty{0.9}{V} हो जाती है। इंजन नियंत्रण वोल्टता कम होने पर ईंधन जोड़ता है और
अधिक होने पर घटाता है, जिससे मिश्रण उस रससमीकरणमितीय बिंदु पर बना रहता है जहाँ
त्रिमार्गी उत्प्रेरक CO, हाइड्रोकार्बनों और नाइट्रोजन ऑक्साइडों को एक साथ रूपांतरित करता है।

\textbf{25.} \qty{700}{\celsius} पर समृद्ध निकास में संवेदक लगभग \qty{0.93}{V} देता है।
\end{solution}
